\documentclass[11pt]{amsart}
\usepackage[margin=1.1in]{geometry}
\usepackage{amsmath,amssymb,amsthm}
\usepackage{enumitem}
\usepackage{url}
\usepackage{hyperref}
\hypersetup{colorlinks=true,linkcolor=blue,citecolor=blue,urlcolor=blue}

\newtheorem{theorem}{Theorem}[section]
\newtheorem{proposition}[theorem]{Proposition}
\newtheorem{lemma}[theorem]{Lemma}
\newtheorem{corollary}[theorem]{Corollary}
\theoremstyle{definition}
\newtheorem{definition}[theorem]{Definition}
\newtheorem{remark}[theorem]{Remark}
\newtheorem{example}[theorem]{Example}

\newcommand{\N}{\mathbb{N}}
\newcommand{\Z}{\mathbb{Z}}
\newcommand{\Q}{\mathbb{Q}}
\newcommand{\R}{\mathbb{R}}
\newcommand{\K}{K}
\newcommand{\pair}[2]{\langle #1,#2\rangle}
\newcommand{\lenbits}[1]{\ell(#1)}
\newcommand{\lean}[1]{\path{#1}}
\newcommand{\names}{\mathcal{S}}
\newcommand{\budget}{b}
\DeclareMathOperator{\dist}{dist}

\title[Mutual information at split precisions]{Mutual information at split precisions:\\ one-sided transfer, symmetry of information, and a Lean formalization}
\author{Emma Dinowitz}
\date{November 7, 2026}

\begin{document}

\begin{abstract}
For points $x\in\R^n$, $y\in\R^t$ Case and Lutz define the mutual information $I_r(x:y)$ of $x$ and $y$ at a common precision $r$.  In applications to smooth dynamics the two points are observed at precisions $r,s$ tied by a conversion rate rather than equal, and we define $I_{r,s}(x:y)$ accordingly.  We prove that $I_{r,s}(x:y)$ equals the mutual information of $K$-minimizers of the two balls up to $O(\log r)+O(\log s)$, and that symmetry of information holds at split precisions,
$I_{r,s}(x:y)=K_r(x)+K_s(y)-K_{r,s}(x,y)+O(\log r)+O(\log s)$, with an error that splits additively between the precisions and never contains a joint term.  The statements hold for tuples of points with error $\sum_i O(\log r_i)+\sum_j O(\log s_j)$, for points of any presented metric space with an effective net structure, and relative to any oracle.  The proofs isolate a single one-sided transfer principle and a small abstract interface (``located families'') under which the whole argument is pure string bookkeeping.  Along the way we observe that the Case--Lutz lower bound on mutual information through a common string is sensitive to the orientation of the asymmetric convention $K(q)-K(q\mid p)$: with the arguments swapped it fails, while under the symmetric convention it holds with constant error and no $K(K(\cdot))$ terms; and we give a direct proof of the Levin--G\'acs direction of symmetry of information for prefix complexity relative to an arbitrary oracle.  Every result is formalized in Lean~4 in the author's Kolmogorov-complexity library; the formal statements are indicated throughout.
\end{abstract}

\maketitle

\section{Introduction}

Case and Lutz \cite{CaseLutz} introduced the mutual information $I_r(x:y)$ of two Euclidean points at a precision $r$ and showed that $I_r(x:y)=K_r(x)+K_r(y)-K_r(x,y)+O(\log r)$, where $K_r(x)$ is the least prefix complexity of a rational point within $2^{-r}$ of $x$.  Lutz and Stull \cite{LutzStull} introduced a two-precision conditional complexity $K_{r,s}(x\mid y)$, but only as a stepping stone: they show that raising $s$ to $r$ costs $O(\log r)$ and then work at a single precision.  In both settings the two scales never carry independent meaning.

The present note is motivated by anisotropic expansion in smooth dynamics, where forward and backward orbit data determine a point at precisions related by Lyapunov potentials rather than equal.  We therefore keep the precisions $r$ and $s$ of $x$ and $y$ independent throughout, and ask what survives.  The answer is: everything, with an error $O(\log r)+O(\log s)$ in which the two scales never interact.  The structural reason is that every transfer step one needs is \emph{one-sided}: it constrains only the object in the conditioning slot of a conditional complexity, while the object being described ranges freely over all strings.  Because no statement ever mentions two precisions at once, no joint error term can arise.

Two further points emerged while formalizing the argument, and we record them because they change the way the results should be stated.

\begin{enumerate}[label=(\alph*)]
\item The inequality
\begin{equation}\label{eq:caselutz}
I(u:v)\;\ge\;K(z)-K(z\mid u,K(u))-K(z\mid v,K(v))-K(K(u))-K(K(z))-O(1),
\end{equation}
which one uses to pass from a rational chosen to minimize mutual information to a $K$-minimizer, is Corollary~4.7 of \cite{CaseLutz} and is correct in their convention $I(p:q)=K(q)-K(q\mid p)$, in which the \emph{second} argument is described (Proposition~\ref{prop:cl}).  It is \emph{false} if the convention is read as $I(u:v)=K(u)-K(u\mid v)$ with the variables of \eqref{eq:caselutz} kept in place (Remark~\ref{rem:convention}), a slip that is easy to make and that the author made in an earlier draft.  Under the symmetric convention $I(u:v)=K(u)+K(v)-K(\pair uv)$ it holds with constant error and without the $K(K(\cdot))$ terms (Proposition~\ref{prop:mi}(iii)), and in fact the proof of the comparison theorem does not need it at all: symmetry of information installs the minimizer in the charged slot directly.
\item Since the intended application is relative to an oracle $A$, the string-level facts are needed for the prefix complexity $K^A$ of an optimal $A$-prefix machine.  The coding direction of symmetry of information relativizes verbatim, but the Levin--G\'acs direction does not follow from its unrelativized form.  We give a self-contained proof for arbitrary $A$ (Theorem~\ref{thm:soi}).
\end{enumerate}

All results below are formalized in Lean~4 \cite{Lean} on top of Mathlib, in the author's library of algorithmic information theory, in the modules \lean{Oracle/PrefixComplexity/} and \lean{PointToSet/SplitPrecision/} (Section~\ref{sec:formal}).  Formal names are given in brackets after each statement.  Every formal theorem depends only on the axioms \lean{propext}, \lean{Classical.choice} and \lean{Quot.sound}.

\section{Preliminaries}\label{sec:prelim}

\subsection{Oracle prefix complexity}
Strings are finite binary words.  Fix an oracle $A\subseteq\N$ and an optimal prefix machine $U=U^A$ relative to $A$: a partial function $U(p,y)$ on pairs of strings, computable with oracle $A$, whose halting set in every context $y$ is prefix-free, and which is optimal among all such machines up to an additive constant.  Such machines exist, and the resulting \emph{conditional prefix complexity}
\[K(x\mid y)=K^A(x\mid y)=\min\{|p|: U(p,y)=x\},\qquad K(x)=K(x\mid\varepsilon),\]
is finite for all $x,y$ and well defined up to $O(1)$ \cite[Ch.~4]{SUV}.  Throughout, $O(1)$ denotes a constant depending only on $A$ (and on the finitely many computable maps mentioned in a statement); $O(\log r)$ denotes $c\,(\lenbits r+1)$ for such a constant $c$, where $\lenbits r$ is the number of binary digits of $r$.  We write $\pair xy$ for the self-delimiting pair code of $x$ and $y$, and $(x,K(x))$ for the context $\pair{x}{\overline{K(x)}}$ in which $K(x)$ is encoded in unary; we use the same conventions for natural numbers through their binary codes.

\begin{proposition}[conditional calculus]\label{prop:calc}
For all strings $x,y,z$ and every $A$-computable total or partial map $f$:
\begin{enumerate}[label=(\roman*)]
\item (chain rule) $K(x\mid z)\le K(y\mid z)+K(x\mid y)+O(1)$ and $K(x\mid z)\le K(y\mid z)+K(x\mid \pair yz)+O(1)$ \ [\lean{oraclePrefixComplexityNat_triangle}, \lean{_triangle_pair}];
\item (coding direction of symmetry) $K(\pair yx\mid z)\le K(y\mid z)+K(x\mid y,K(y\mid z))+O(1)$ \ [\lean{oraclePrefixComplexityNat_pair_le_cond_context}];
\item (concatenation) $K(\pair yx\mid z)\le K(y\mid z)+K(x\mid z)+O(1)$ \ [\lean{oraclePrefixComplexityNat_pair_le_add}];
\item (recoding) $K(x\mid y)\le K(x\mid f(y))+O(1)$ and $K(f(x)\mid y)\le K(x\mid y)+O(1)$ \ [\lean{_cond_le_of_partrecIn}, \lean{_map_le_of_partrecIn}];
\item $K(x\mid y)\le K(x\mid y,k)+K(\overline k)+O(1)$ for every $k\in\N$, and $K(\overline k)\le 3\lenbits k+O(1)$ \ [\lean{_le_context_add}, \lean{oracleComplexity_natCode_le}].
\end{enumerate}
\end{proposition}

\begin{proof}
(i)--(iii) are instances of one machine.  Given $A$-computable maps $g$ and $h$, let $D(p_1p_2,z)=h(z,y,x)$ where $U(p_1,z)=y$ and $U(p_2,g(z,y,|p_1|))=x$.  Prefix-freeness of the domain of $U$ in each context makes the parse $p_1p_2$ unique, so $D$ is again a prefix machine, and it is $A$-computable by dovetailing over the enumerable set of halting pairs $(p_1,p_2)$.  Optimality of $U$ gives $K(h(z,y,x)\mid z)\le K(y\mid z)+K(x\mid g(z,y,K(y\mid z)))+O(1)$ [\lean{oraclePrefixComplexityNat_twoStage_le}]; take $g=y$, $g=\pair yz$, $g=(y,|p_1|)$ and $g=z$ respectively.  (iv) composes $U$ with $f$ on the context or the output.  (v) follows from (iv) and (iii) applied to $\overline k$, and from the binary code of $k$ having $\lenbits k$ bits.
\end{proof}

\subsection{Symmetry of information relative to an oracle}

\begin{theorem}[symmetry of information for $K^A$]\label{thm:soi}
For every oracle $A$ and all strings $x,y$,
\[K(\pair xy)=K(x)+K(y\mid x,K(x))+O(1).\]
\textup{[\lean{oracleComplexity_pair_le_cond_context}, \lean{oracleSymmetryLower}]}
\end{theorem}

The inequality $\le$ is Proposition~\ref{prop:calc}(ii) with $z=\varepsilon$.  For $\ge$ we cannot invoke the unrelativized theorem, whose proof is about computable machines; for oracles with computable characteristic function the transport is immediate [\lean{OracleSymmetryLower.of_partrec}], but in general a direct argument is needed.

\begin{proof}[Proof of $\ge$]
For a string $x$ let $\mu(x)$ be the total weight $\sum 2^{-|p|}$ of the programs $p$ with $U(p,\varepsilon)=\pair x{y}$ for some $y$.  Since the domain of $U$ is prefix-free, $\sum_x\mu(x)\le 1$, and since distinct outputs come from distinct programs, $\sum_y 2^{-K(\pair xy)}\le\mu(x)$.

\emph{Marginal coding bound} [\lean{oracle_pairMarginal_coding_bound}]: $\mu(x)\le 2^{c}\,2^{-K(x)}$ for a constant $c$.  The set $R=\{(x,n):2^{-n}<\mu(x)\}$ is $A$-enumerable, because $\mu(x)$ is the supremum of the rational partial sums over the enumerable halting triples $(p,x,y)$.  Enumerate $R$ without repetition and issue, for the pair $(x,n)$, a Kraft--Chaitin request for $x$ of length $n+1$.  The total mass is at most $\sum_x\sum_{n:2^{-n}<\mu(x)}2^{-(n+1)}\le\sum_x\mu(x)\le 1$, so by the relative Kraft--Chaitin theorem $K(x)\le n+1+O(1)$ whenever $2^{-n}<\mu(x)$; taking $n$ least gives $\mu(x)\le 2^{-n+1}\le 2^{c}2^{-K(x)}$.

\emph{Conditional requests} [\lean{oracle_kraftChaitin_cond}]: the Kraft--Chaitin construction is uniform in a context string, so an $A$-computable family of request streams indexed by contexts is realized by a single $A$-prefix machine, exactly at the requested lengths, in every context whose stream has mass at most one.  Given the context $(x,k)$, enumerate the halting triples $(p,x,y)$ and request $y$ with length $|p|-k+c+c_2$, where $c_2$ is the constant of $K(x)\le K(\pair xy)+c_2$ (so no truncation occurs when $k=K(x)$).  At $k=K(x)$ the mass is $2^{k-c-c_2}\mu(x)\le 2^{-c_2}\le 1$, hence $K(y\mid x,K(x))\le K(\pair xy)-K(x)+O(1)$.
\end{proof}

\subsection{Mutual information of strings}

\begin{definition}
The \emph{mutual information} of two strings is $I(u:v)=K(u)+K(v)-K(\pair uv)\in\Z$ \ [\lean{oracleMutualInfo}].
\end{definition}

\begin{proposition}\label{prop:mi}
For all strings $u,v,z$:
\begin{enumerate}[label=(\roman*)]
\item $I(u:v)\ge -O(1)$ and $I(u:v)=I(v:u)+O(1)$;
\item $I(u:v)=K(v)-K(v\mid u,K(u))+O(1)$ \ [\lean{oracleMutualInfo_ge_cond}, \lean{oracleMutualInfo_le_cond}];
\item \label{it:common} $I(u:v)\ge K(z)-K(z\mid u,K(u))-K(z\mid v,K(v))-O(1)$ \ [\lean{oracleMutualInfo_ge_common}].
\end{enumerate}
\end{proposition}

\begin{proof}
(i) is Proposition~\ref{prop:calc}(iii),(iv); (ii) is Theorem~\ref{thm:soi}.  For (iii), $K(\pair uv)\le K(z)+K(\pair uv\mid z,K(z))+O(1)\le K(z)+K(u\mid z,K(z))+K(v\mid z,K(z))+O(1)$ by Proposition~\ref{prop:calc}(ii),(iii), while $K(u)\ge K(z)+K(u\mid z,K(z))-K(z\mid u,K(u))-O(1)$ and likewise for $v$ by both directions of Theorem~\ref{thm:soi}; add.
\end{proof}

\begin{proposition}[the Case--Lutz convention]\label{prop:cl}
Let $I_{\mathrm{CL}}(u:v)=K(v)-K(v\mid u)$.  Then $I_{\mathrm{CL}}(u:v)\le I(u:v)+O(1)$, $I(u:v)\le I_{\mathrm{CL}}(u:v)+K(\overline{K(u)})+O(1)$, and for all $x,y,z$
\[I_{\mathrm{CL}}(x:y)\;\ge\;K(z)-K(z\mid x,K(x))-K(z\mid y,K(y))-K(\overline{K(x)})-K(\overline{K(z)})-O(1),\]
which is Corollary~4.7 of \cite{CaseLutz}. \ \textup{[\lean{oracleMutualInfoCL_le}, \lean{oracleMutualInfo_le_CL_add}, \lean{oracleMutualInfoCL_ge_common}]}
\end{proposition}

\begin{proof}
The comparisons are Proposition~\ref{prop:mi}(ii) with Proposition~\ref{prop:calc}(iv),(v).  For the inequality, both directions of Theorem~\ref{thm:soi} give $K(y)\ge K(z)+K(y\mid z,K(z))-K(z\mid y,K(y))-O(1)$, and the chain rule with Proposition~\ref{prop:calc}(v) gives $K(y\mid x)\le K(z\mid x,K(x))+K(\overline{K(x)})+K(y\mid z,K(z))+K(\overline{K(z)})+O(1)$; subtract.  The two $K(K(\cdot))$ terms are exactly the cost of supplying the complexity values of $x$ and $z$ to the chain through $z$.
\end{proof}

\begin{remark}[the orientation matters]\label{rem:convention}
With the arguments swapped, $I'(u:v)=K(u)-K(u\mid v)$, the same inequality \eqref{eq:caselutz} is false.  Take $z=u$ and $u=\overline{K(v)}$.  Then $K(z\mid u,K(u))=O(1)$, $K(z\mid v,K(v))=O(1)$, and $K(K(u))$, $K(K(z))$ are $O(\log\log|v|)$, so \eqref{eq:caselutz} would give $K(K(v)\mid v)=K(u\mid v)+O(1)\le O(\log\log |v|)$ for all $v$, contradicting G\'acs' theorem that $K(K(v)\mid v)$ exceeds $\log|v|-O(\log\log|v|)$ infinitely often \cite{Gacs}.  Since $I'(u:v)-I(u:v)=K(u\mid v,K(v))-K(u\mid v)+O(1)\in[-K(K(v)\mid v)-O(1),\,O(1)]$ and $I_{\mathrm{CL}}$ differs from $I$ by at most $K(\overline{K(u)})+O(1)$, all three conventions differ by $O(\log r)$ at $K$-minimizers of balls of radius $2^{-r}$ (Proposition~\ref{prop:calc}(v) with $K(\cdot)=O(r)$), so all statements about points below hold in each of them; but only the symmetric one makes the string facts exact and orientation-free.
\end{remark}

\section{Located families and one-sided transfer}\label{sec:located}

The objects of interest, points and tuples of points, are known only through \emph{names}: strings ranging over a set determined by the object and a precision.  We abstract exactly what the argument uses.

\begin{definition}\label{def:located}
Let $\iota$ be a set of \emph{precisions}.  A \emph{located family} $L=(\names,\budget)$ over $\iota$ consists of nonempty sets of strings $\names_r$ and budgets $\budget_r\in\N$, $r\in\iota$, such that for a constant $c$ and all $r$:
\begin{enumerate}[label=(L\arabic*)]
\item\label{L1} (transfer) $K(p\mid q)\le c\,(\budget_r+1)$ for every $q\in\names_r$ and every $K$-minimizer $p$ of $\names_r$;
\item\label{L2} (growth) $\lenbits{K(p)}\le c\,(\budget_r+1)$ for every $K$-minimizer $p$ of $\names_r$.
\end{enumerate}
Here a $K$-minimizer of $\names_r$ is an element of least complexity; one exists since complexities are natural numbers.  We write $K_L(r)$ for that least complexity.  \textup{[\lean{LocatedFamily}, \lean{IsMinName}, \lean{LocatedFamily.complexity}]}
\end{definition}

\begin{lemma}[transfer lemmas]\label{lem:transfer}
Let $L$ be located over $\iota$.  For all $r\in\iota$, $q\in\names_r$, $K$-minimizers $p$ of $\names_r$ and all strings $z,w$:
\begin{enumerate}[label=(\roman*)]
\item $K(z\mid \pair qw)\le K(z\mid\pair pw)+O(\budget_r)$ \ [\lean{LocatedFamily.transfer_cond}];
\item $K(z\mid q,K(q))\le K(z\mid p,K(p))+O(\budget_r)$ \ [\lean{LocatedFamily.transfer_context}];
\item $K(\overline{K(p)})\le O(\budget_r)$ \ [\lean{LocatedFamily.natCode_complexity_le}].
\end{enumerate}
\end{lemma}

\begin{proof}
(iii) is \ref{L2} and Proposition~\ref{prop:calc}(v).  For (i), $K(z\mid\pair qw)\le K(\pair pw\mid\pair qw)+K(z\mid\pair pw)+O(1)$ and $K(\pair pw\mid \pair qw)\le K(p\mid q)+K(w\mid\pair qw)+O(1)\le O(\budget_r)$ by \ref{L1}.  For (ii), the same chain with $K((p,K(p))\mid (q,K(q)))\le K(p\mid q)+K(\overline{K(p)})+O(1)\le O(\budget_r)$.
\end{proof}

Each lemma constrains only the string in the conditioning slot; $z$ and $w$ are arbitrary.  This is what allows two located families with unrelated precisions to be treated one at a time.

\begin{proposition}[products]\label{prop:prod}
Let $L_1$, $L_2$ be located over $\iota_1$, $\iota_2$.  Then $L_1\times L_2$, with names $\{\pair{q_1}{q_2}: q_i\in\names^{(i)}_{r_i}\}$ and budgets $\budget^{(1)}_{r_1}+\budget^{(2)}_{r_2}$, is located over $\iota_1\times\iota_2$.  Moreover, if $p_1,p_2$ are $K$-minimizers of $\names^{(1)}_{r_1}$, $\names^{(2)}_{r_2}$, then
\[K_{L_1\times L_2}(r_1,r_2)\le K(\pair{p_1}{p_2})\le K_{L_1\times L_2}(r_1,r_2)+O(\budget^{(1)}_{r_1}+\budget^{(2)}_{r_2}).\]
\textup{[\lean{LocatedFamily.prod}, \lean{LocatedFamily.complexity_pair_min_le}]}
\end{proposition}

\begin{proof}
A joint minimizer $p=\pair{u_1}{u_2}$ of the product is in general not a pair of marginal minimizers, nor conversely, so both notions are kept.  Let $p^*=\pair{p_1}{p_2}$.  By \ref{L1} for the factors and concatenation, $K(p^*\mid\pair{q_1}{q_2})$ and $K(p^*\mid p)$ are $O(\budget^{(1)}+\budget^{(2)})$.  By Theorem~\ref{thm:soi}, $K(p\mid p^*,K(p^*))\le K(\pair{p^*}{p})-K(p^*)+O(1)\le K(p)+K(p^*\mid p)-K(p^*)+O(1)\le K(p^*\mid p)+O(1)$, using $K(p)\le K(p^*)$.  Then $K(p\mid\pair{q_1}{q_2})\le K(p\mid p^*)+K(p^*\mid\pair{q_1}{q_2})+O(1)$ and $K(p\mid p^*)\le K(p\mid p^*,K(p^*))+K(\overline{K(p^*)})+O(1)$, which is \ref{L1}; \ref{L2} follows from $K(p)\le K(p^*)\le K(p_1)+K(p_2)+O(1)$.  The displayed inequality is $K(p^*)\le K(p)+K(p^*\mid p)+O(1)$.
\end{proof}

Iterating, a tuple $(L_1,\dots,L_k)$ of located families is located over $\prod\iota_i$, with names the list codes of tuples of names and budget $\sum_i\budget^{(i)}_{r_i}$ [\lean{LocatedFamily.pi}, \lean{mem_pi_names}, \lean{pi_budget}].  This is the coordinatewise iteration of the transfer lemmas, and the proposition is the comparison of joint and marginal minimizers.

\section{Mutual information at split precisions}\label{sec:split}

\begin{definition}
For located families $L_1,L_2$ over $\iota_1,\iota_2$ and $(r,s)\in\iota_1\times\iota_2$ put
\[I_{r,s}(L_1:L_2)=\min\{I(q_1:q_2): q_i\in\names^{(i)}\},\qquad J_{r,s}(L_1:L_2)=\min\{I(p_1:p_2): p_i\text{ a $K$-minimizer of }\names^{(i)}\},\]
the minima over $\names^{(1)}_r\times\names^{(2)}_s$ being attained since $I\ge -O(1)$. \textup{[\lean{LocatedFamily.splitInfo}, \lean{LocatedFamily.minInfo}]}
\end{definition}

Minimizers are names, so $I_{r,s}\le J_{r,s}$.  The two precisions do not interact in the definitions, and the following theorems show that they do not interact in the errors either.

\begin{theorem}[comparison]\label{thm:comparison}
$J_{r,s}(L_1:L_2)\le I_{r,s}(L_1:L_2)+O(\budget^{(1)}_r+\budget^{(2)}_s)$. \ \textup{[\lean{LocatedFamily.minInfo_le_splitInfo_add}]}
\end{theorem}

\begin{proof}
Let $q_1,q_2$ attain $I_{r,s}$ and let $p_1,p_2$ be arbitrary minimizers.  Every step below is charged to one precision.

\emph{Step 1 (charged to $r$).}  By Proposition~\ref{prop:mi}(ii), $I(q_1:q_2)\ge K(q_2)-K(q_2\mid q_1,K(q_1))-O(1)$.  By the chain rule and Lemma~\ref{lem:transfer}, $K(q_2\mid q_1,K(q_1))\le K(p_1\mid q_1)+K(q_2\mid p_1)+O(1)\le K(q_2\mid p_1,K(p_1))+O(\budget^{(1)}_r)$, using Lemma~\ref{lem:transfer}(iii) to add $K(p_1)$ to the context.  Hence $I(q_1:q_2)\ge K(q_2)-K(q_2\mid p_1,K(p_1))-O(\budget^{(1)}_r)\ge I(p_1:q_2)-O(\budget^{(1)}_r)$, again by Proposition~\ref{prop:mi}(ii).  The $y$-side is unconstrained, so $s$ does not enter.

\emph{Step 2 (charged to $s$).}  Symmetrically, $I(p_1:q_2)=I(q_2:p_1)+O(1)\ge K(p_1)-K(p_1\mid q_2,K(q_2))-O(1)\ge K(p_1)-K(p_1\mid p_2,K(p_2))-O(\budget^{(2)}_s)$ by the decorated transfer lemma~\ref{lem:transfer}(ii), and the last expression is $I(p_2:p_1)+O(1)=I(p_1:p_2)+O(1)$.

Since $p_1,p_2$ were arbitrary, $J_{r,s}\le I(p_1:p_2)\le I_{r,s}+O(\budget^{(1)}_r+\budget^{(2)}_s)$.
\end{proof}

The proof is the four-step argument of the original appendix with its Step~2 removed: there the minimizer was installed in the charged slot of \eqref{eq:caselutz}, which is exactly where the orientation of Remark~\ref{rem:convention} bites.  Under the symmetric convention Proposition~\ref{prop:mi}(ii) does the installation.

\begin{theorem}[symmetry of information at split precisions]\label{thm:symmetry}
\[I_{r,s}(L_1:L_2)=K_{L_1}(r)+K_{L_2}(s)-K_{L_1\times L_2}(r,s)+O(\budget^{(1)}_r+\budget^{(2)}_s).\]
\textup{[\lean{LocatedFamily.splitInfo_symmetry}]}
\end{theorem}

\begin{proof}
Let $p_1,p_2$ attain $J_{r,s}$.  Then $J_{r,s}=K(p_1)+K(p_2)-K(\pair{p_1}{p_2})=K_{L_1}(r)+K_{L_2}(s)-K(\pair{p_1}{p_2})$ and $K(\pair{p_1}{p_2})=K_{L_1\times L_2}(r,s)+O(\budget^{(1)}+\budget^{(2)})$ by Proposition~\ref{prop:prod}; conclude with Theorem~\ref{thm:comparison}.
\end{proof}

\begin{corollary}\label{cor:sym}
$I_{r,s}(L_1:L_2)=I_{s,r}(L_2:L_1)+O(1)$. \ \textup{[\lean{LocatedFamily.splitInfo_comm_le}]}
\end{corollary}

\begin{corollary}[normalized forms]\label{cor:normalized}
Let $(r_k,s_k)$ be precisions and $N_k\to\infty$ a scale with $\budget^{(1)}_{r_k}/N_k\to0$ and $\budget^{(2)}_{s_k}/N_k\to0$.  Then
\[\frac{J_{r_k,s_k}-I_{r_k,s_k}}{N_k}\to0\qquad\text{and}\qquad\frac{I_{r_k,s_k}-\bigl(K_{L_1}(r_k)+K_{L_2}(s_k)-K_{L_1\times L_2}(r_k,s_k)\bigr)}{N_k}\to0.\]
\textup{[\lean{tendsto_minInfo_sub_splitInfo_div}, \lean{tendsto_splitInfo_sub_symmetry_div}]}
\end{corollary}

Consequently the lower and upper limits of the normalized mutual information agree with those of the normalized symmetry expression.  Some such condition is close to necessary: if $s$ is held fixed while $r\to\infty$, Theorem~\ref{thm:comparison} still holds but both sides are bounded by $O(s)$, so the normalized statement carries no information.

\section{Points, tuples and Euclidean space}\label{sec:points}

\subsection{Presented spaces}
A \emph{presented} metric space is a complete metric space $X$ with a fixed dense sequence $(\xi_i)_{i\in\N}$ of \emph{ideal points} \cite{DinowitzPTS}.  The complexity of a point at precision $r$ is
\[K_r(x)=\min\{K(\bar i): \dist(\xi_i,x)<2^{-r}\},\]
where $\bar i$ is the binary code of $i$; its minimizers are the $K$-minimizers of the ball.  For $X=\R^n$ with the sup metric and $(\xi_i)$ an enumeration of $\Q^n$ this is the usual $K_r(x)$ of \cite{CaseLutz,LutzStull}.  The names of $x$ at precision $r$ are the codes of the ideal points of $B(x,2^{-r})$ [\lean{ballNames}], and the least complexity of a name is $K_r(x)$ [\lean{minComplexity_ballNames}].

\begin{definition}\label{def:nets}
A presented space has \emph{precision nets} relative to $A$ if there are an $A$-computable function $(r,i,j)\mapsto\mathrm{cand}(r,i,j)\in\N$ and a bound $D$ such that every ball $B(x,2^{-r})$ contains a \emph{canonical} ideal point $\xi_d$ with $d=\mathrm{cand}(r,i,j)$ for some $j<D$, for every ideal point $\xi_i$ of the ball.  A point $x$ has \emph{linear growth} if $K_r(x)\le C(r+1)$ for some $C$. \textup{[\lean{PrecisionNets}, \lean{HasLinearGrowth}]}
\end{definition}

Linear growth is the only place the ambient geometry enters the estimates (and it enters only through the number of bits of $K_r(x)$); any bound $K_r(x)\le Cr+O(\log r)$ would do.

\begin{theorem}[geometric transfer]\label{thm:geomtransfer}
Let $X$ have precision nets relative to $A$ and let $x$ have linear growth.  Then $x$ is a located family over $\N$ with names the codes of the ideal points of $B(x,2^{-r})$ and budget $\lenbits r$: for every ideal point $\xi_i$ of $B(x,2^{-r})$ and every $K$-minimizer $\xi_p$ of the ball, $K(\bar p\mid\bar i)\le O(\log r)$.
\textup{[\lean{transfer_ballNames}, \lean{pointFamily}]}
\end{theorem}

\begin{proof}
Let $\xi_d$ be the canonical point of $B(x,2^{-r})$.  From any ideal point $\xi_i$ of the ball, $d$ is recovered from $r$ and an index $j<D$, so $K(\bar d\mid\bar i)\le K(\bar r)+K(\bar j)+O(1)=O(\log r)$ [\lean{PrecisionNets.canonical_cond_le}]; in particular $K(\bar d)\le K(\bar p)+O(\log r)\le C'(r+1)$, so $K(\overline{K(\bar d)})=O(\log r)$.  Since $\xi_d$ lies in the ball, $K(\bar p)\le K(\bar d)$, and Theorem~\ref{thm:soi} gives $K(\bar p\mid \bar d,K(\bar d))\le K(\pair{\bar d}{\bar p})-K(\bar d)+O(1)\le K(\bar p)+K(\bar d\mid\bar p)-K(\bar d)+O(1)\le O(\log r)$.  Hence $K(\bar p\mid\bar i)\le K(\bar d\mid\bar i)+K(\bar p\mid\bar d)+O(1)\le O(\log r)+K(\bar p\mid\bar d,K(\bar d))+K(\overline{K(\bar d)})+O(1)=O(\log r)$.  The growth condition \ref{L2} is linear growth.
\end{proof}

This is the content of the transfer lemmas whose proofs were deferred in the original appendix: the canonical point replaces the ``pivot through the dyadic truncation'', and symmetry of information replaces the count of almost-realizers.  No counting of rationals of bounded complexity is needed.

\subsection{The theorems for points}
Let $X,Y$ be presented spaces with precision nets relative to $A$, and let $x\in X$, $y\in Y$ have linear growth.  Define
\[I_{r,s}(x:y)=\min\{I(\bar i:\bar j):\dist(\xi_i,x)<2^{-r},\ \dist(\eta_j,y)<2^{-s}\},\]
$J_{r,s}(x:y)$ the same minimum over $K$-minimizers of the two balls, and $K_{r,s}(x,y)=\min K(\pair{\bar i}{\bar j})$ over the same pairs.  These definitions involve no hypotheses [\lean{splitMutualInfo}, \lean{minimizerMutualInfo}, \lean{jointPrecisionComplexity}]; for $r=s$ they are the quantities of \cite{CaseLutz}.

\begin{theorem}\label{thm:points}
For all $r,s\in\N$:
\begin{enumerate}[label=(\roman*)]
\item $I_{r,s}(x:y)\le J_{r,s}(x:y)\le I_{r,s}(x:y)+O(\log r)+O(\log s)$ \ [\lean{minimizerMutualInfo_le_splitMutualInfo_add}];
\item $I_{r,s}(x:y)=K_r(x)+K_s(y)-K_{r,s}(x,y)+O(\log r)+O(\log s)$ \ [\lean{splitMutualInfo_symmetry}];
\item $I_{r,s}(x:y)=I_{s,r}(y:x)+O(1)$ \ [\lean{splitMutualInfo_comm_le}].
\end{enumerate}
If $(r_k,s_k)$ is log-comparably divergent, i.e.\ $r_k,s_k\to\infty$, $\log s_k/r_k\to0$ and $\log r_k/s_k\to0$ --- in particular if it is comparably divergent, $0<\liminf r_k/s_k\le\limsup r_k/s_k<\infty$ --- then $(J_{r_k,s_k}-I_{r_k,s_k})/r_k\to0$ and the normalized mutual information and symmetry expression have the same limiting behaviour; the same holds with $s_k$ in the denominator.
\textup{[\lean{tendsto_minimizerMutualInfo_sub_splitMutualInfo_div}, \lean{tendsto_splitMutualInfo_sub_symmetry_div}]}
\end{theorem}

\begin{proof}
Theorems~\ref{thm:comparison}, \ref{thm:symmetry}, Corollaries~\ref{cor:sym}, \ref{cor:normalized} applied to the located families of Theorem~\ref{thm:geomtransfer}, whose product has least complexity $K_{r,s}(x,y)$ and budget $\lenbits r+\lenbits s$; and $\lenbits{s_k}/r_k\to0$ because $\lenbits s\le\log_2 s+1$.
\end{proof}

\subsection{Tuples}
For $k$ points $x_i\in X_i$ at precisions $\bar r=(r_1,\dots,r_k)$ and $m$ points $y_j\in Y_j$ at $\bar s$, the names of $\bar x$ at $\bar r$ are the list codes of tuples of ideal points of the balls $B(x_i,2^{-r_i})$, $K_{\bar r}(\bar x)$ is the least complexity of such a name, and $I_{\bar r,\bar s}(\bar x:\bar y)$, $J_{\bar r,\bar s}$, $K_{\bar r,\bar s}(\bar x,\bar y)$ are defined as before [\lean{tupleNames}, \lean{tuplePrecisionComplexity}, \lean{tupleSplitMutualInfo}, \lean{tupleJointComplexity}].  Each $x_i$ is located with budget $\lenbits{r_i}$, so by the product construction $\bar x$ is located with budget $\sum_i\lenbits{r_i}$, and Section~\ref{sec:split} applies verbatim.

\begin{theorem}\label{thm:tuples}
$J_{\bar r,\bar s}(\bar x:\bar y)\le I_{\bar r,\bar s}(\bar x:\bar y)+\sum_iO(\log r_i)+\sum_jO(\log s_j)$, and
\[I_{\bar r,\bar s}(\bar x:\bar y)=K_{\bar r}(\bar x)+K_{\bar s}(\bar y)-K_{\bar r,\bar s}(\bar x,\bar y)+\sum_iO(\log r_i)+\sum_jO(\log s_j).\]
\textup{[\lean{tupleMinimizerMutualInfo_le_tupleSplitMutualInfo_add}, \lean{tupleSplitMutualInfo_symmetry}]}
\end{theorem}

No step of any proof holds an $r_i$-constrained and an $s_j$-constrained object in constrained position simultaneously, nor two $r_i$-constrained objects; this is why the error is a sum over all $k+m$ precisions and never a joint term.

\subsection{Euclidean space}
\begin{proposition}\label{prop:euclid}
$\R^n$ with the sup metric and the rational points as ideal points is a presented space with an exact computable distance table, with precision nets relative to every oracle, and every point has linear growth.  Consequently Theorems~\ref{thm:points} and \ref{thm:tuples} hold for $x\in\R^n$, $y\in\R^t$ and every oracle $A$ without further hypotheses.
\textup{[\lean{euclidPresentation}, \lean{euclidNets}, \lean{hasLinearGrowth_euclid}, \lean{euclid_minimizerMutualInfo_le_splitMutualInfo_add}, \lean{euclid_splitMutualInfo_symmetry}]}
\end{proposition}

\begin{proof}
The canonical point of $B(x,2^{-r})$ is the dyadic truncation $d=(\lfloor 2^{r+2}x_i\rfloor/2^{r+2})_i$, within $2^{-r-2}$ of $x$.  If $q$ is a rational point of the ball then $|\lfloor2^{r+2}x_i\rfloor-\lfloor2^{r+2}q_i\rfloor|\le4$ for each $i$, so $d$ is one of the $9^n$ points $\bigl((\lfloor2^{r+2}q_i\rfloor+e_i-4)/2^{r+2}\bigr)_i$, $e\in\{0,\dots,8\}^n$, which are computable from $(r,q,e)$.  For linear growth, the code of $d$ is computable from the list $(r,\lfloor2^{r+2}x_1\rfloor,\dots,\lfloor2^{r+2}x_n\rfloor)$, whose entries have $O(r+\log(|x_i|+1))$ bits, so $K_r(x)\le K(\bar d)\le n r+O_x(\log r)\le C_x(r+1)$.  The Euclidean metric would do equally well; the sup metric is used because the distance between rational points is then rational.
\end{proof}

Any doubling metric space with an $A$-computable family of finite nets also has precision nets, so the results hold there directly; the pull-back along an Assouad embedding is not needed.

\section{The formalization}\label{sec:formal}

The development lives in the author's Lean~4 library of algorithmic information theory (Mathlib \lean{v4.35.0}).  The string layer is in \lean{KolmogorovMathlib/Oracle/PrefixComplexity/}: \lean{TwoStage.lean} (the two-stage machine and Proposition~\ref{prop:calc}), \lean{Conditional.lean} (the conditional calculus), \lean{Information.lean} (mutual information, Proposition~\ref{prop:mi}), \lean{AsymmetricInformation.lean} (Proposition~\ref{prop:cl}) and \lean{SymmetryLower.lean} (Theorem~\ref{thm:soi}), with the context-uniform Kraft--Chaitin theorem \lean{oracle_kraftChaitin_cond} in \lean{Oracle/KraftChaitin.lean}.  The split-precision theory is in \lean{KolmogorovMathlib/PointToSet/SplitPrecision/}: \lean{Located.lean} (Section~\ref{sec:located}), \lean{Information.lean} and \lean{Normalized.lean} (Section~\ref{sec:split}), \lean{Point.lean}, \lean{Theorems.lean}, \lean{Tuples.lean}, \lean{Euclidean.lean} and \lean{EuclideanTheorems.lean} (Section~\ref{sec:points}); about 4\,400 lines in all.  Every statement is relative to an arbitrary oracle $A\subseteq\N$, error terms are explicit existential constants multiplying $\lenbits r+\lenbits s+1$ (or the sum over coordinates), complexities are natural numbers and mutual informations integers, and the normalized statements are limits of real sequences.  All declarations depend only on the axioms \lean{propext}, \lean{Classical.choice} and \lean{Quot.sound}, and the library's audit (no \lean{sorry}, complete documentation, no resource overrides) passes on the new modules.  The blueprint \lean{docs/split_precision_blueprint.md} maps the statements above to the Lean declarations.

Not formalized: a conditional form $K_{\bar r,\bar s}(\bar x\mid\bar y)$ in the style of \cite{LutzStull}, whose convention (maximum or minimum over the $y$-names) is left open here.

\begin{thebibliography}{9}
\bibitem{CaseLutz} A.~Case and J.~H.~Lutz, \emph{Mutual dimension}, ACM Trans.\ Comput.\ Theory \textbf{7} (2015), no.~3, Art.~12.
\bibitem{LutzStull} N.~Lutz and D.~M.~Stull, \emph{Bounding the dimension of points on a line}, Inform.\ and Comput.\ \textbf{275} (2020); \emph{Projection theorems using effective dimension}, MFCS 2018.
\bibitem{SUV} A.~Shen, V.~A.~Uspensky and N.~Vereshchagin, \emph{Kolmogorov complexity and algorithmic randomness}, Math.\ Surveys Monogr.\ 220, Amer.\ Math.\ Soc., 2017.
\bibitem{Gacs} P.~G\'acs, \emph{On the symmetry of algorithmic information}, Soviet Math.\ Dokl.\ \textbf{15} (1974), 1477--1480.
\bibitem{DinowitzPTS} E.~Dinowitz, \emph{Point-to-set principles in dynamical systems}, preprint, arXiv:2607.21976.
\bibitem{Lean} L.~de~Moura and S.~Ullrich, \emph{The Lean 4 theorem prover and programming language}, CADE 28 (2021); The Mathlib Community, \emph{The Lean mathematical library}, CPP 2020.
\end{thebibliography}

\end{document}
